\documentclass{article}
\usepackage{/globals/de}
\usepackage{/globals/anki}
\ankiprefix{gemeinsamer-ausschuss}
\ankitopic{Gemeinsamer Ausschuss}
\ankishow
\begin{document}
\section{Der Gemeinsame Ausschuss} 
Der Gemeinsame Ausschuss wird als \textquote{Notparlament} tätig wenn es zu einem Verteidigungsfall gekommen ist und der Bundestag nicht (schnell genug) zusammentreten kann.

Er kann den Bundestag und den Bundesrat ersetzen.
\ankinote{aufgabe}{Aufgabe}{Notparlament, wenn BT in Verteidigungsfall nicht zusammentreten kann}

Er besteht zu $2/3$ aus Mitgliedern des Bundestages und zu $1/3$ aus Mitgliedern das Bundesrates.
\ankinote{zusammensetzung}{Zusammensetzung}{32 MdBundestag, 16 MdBundesrat}

\subsection{Verteidigungsfall}
Der Verteidigungsfall ist der Fall in dem das Bundesgebiet mit Waffengewalt (dazu zählt auch ein Cyberangriff) angegriffen wird, oder es unmittelbar droht dass dies passiert.

Dieser wird vom Bundestag festgestellt.
\ankinote{verteidigungsfall}{Verteidigungsfall}{Bundesgebiet Waffengewalt}
\end{document}